\newpage

\section{মীযান}

\noindent যে শাস্ত্রে শব্দ গঠন ও রূপান্তরের নিয়ম জানা যায়, তাকে
{\arabicfont \textarabic{علم الصرف}} বলে।
আলোচ্য আরবি শব্দ। উদ্দেশ্য শাব্দিক ভুলে আটকানো।

\vspace{0.35cm}

\noindent {\arabicfont \textarabic{كلمة}} মানে একক, অর্থবোধক শব্দ। তিন প্রকার।

\vspace{0.2cm}

\begingroup
\setlength{\arrayrulewidth}{0.5pt}
\setlength{\tabcolsep}{4pt}
\renewcommand{\arraystretch}{1.7}

\noindent\begin{tabularx}{\textwidth}{|C|C|C|}
  \hline
  \rowcolor{headergreen}
  {\color{white}\textbf{প্রকার}} &
  {\color{white}\textbf{লক্ষণ}} &
  {\color{white}\textbf{যেমন}} \\
  \hline
  \rowcolor{lightgreen}
  {\arabicfont \textarabic{اسم}} &
  কাল নেই &
  {\arabicfont \textarabic{رَجُلٌ}} পুরুষ \\
  \hline
  {\arabicfont \textarabic{فعل}} &
  তিন কালের একটি আছে &
  {\arabicfont \textarabic{ضَرَبَ}} মারল \\
  \hline
  \rowcolor{lightgreen}
  {\arabicfont \textarabic{حرف}} &
  একা অর্থ জানায় না &
  {\arabicfont \textarabic{مِنْ}} থেকে \\
  \hline
\end{tabularx}
\endgroup

\vspace{0.4cm}

\noindent ফেলের তিন মূল অক্ষর {\arabicfont \textarabic{ف ع ل}},
নাম ফা, আইন, লাম।
মাজির শেষে যবর, মুজারের শেষে পেশ।
দ্বিবচনের আলিফ বা বহুবচনের ওয়াও এলে সেই শেষ হরকত বদলায়।
হাল ও ইস্তিকবালের সিগা এক। দুটোর নাম মুজারে।

\vspace{0.3cm}

\noindent মারুফ: কর্তা জ্ঞাত।
মাজহুল: কর্তা অজ্ঞাত।
মাজহুল মাজি {\arabicfont \majhulmadi},
মুজারে {\arabicfont \majhulmud}।

\vspace{0.25cm}

\noindent পুরুষ তিন (গায়েব, হাজির, মুতাকাল্লিম), বচন তিন, লিঙ্গ দুই।
মুতাকাল্লিমে লিঙ্গ আলাদা হয় না। মোট চৌদ্দ সিগা।
মাজি ছয় রকম। পরের ছক প্রথমটি: মুতলাক মারুফ।
